% T11 -- Accuracy by how much history a user has (slide 77).
\begin{table}[t]
\centering\small
\caption[Accuracy by record count]{The 1{,}000 held-out users of \seleven{} cut
into six equal groups by how many records they have. The persistence baseline is
right 17.8\,\% of the time on a user with fewer than 30 records and 59.8\,\% of
the time on a user with more than 120: a 42-point spread, against 0.3\pt\ for
the whole training-size ladder. Length of history, not size of training set, is
the axis of difficulty. 21{,}000-user model, 80\,\% context.}
\label{tab:record-count}
\begin{tabular}{@{}llrrrrr@{}}
\toprule
& \textbf{Records} & \textbf{Users} & \textbf{Predictions} & \textbf{Baseline} & \textbf{Model} & \textbf{Gap} \\
\midrule
R1 & 10--29   & 161 &   695 & 17.8\,\% & 22.1\,\% & $+4.3$ \\
R2 & 30--43   & 168 & 1{,}221 & 37.7\,\% & 42.0\,\% & $+4.3$ \\
R3 & 44--59   & 159 & 1{,}635 & 43.1\,\% & 47.4\,\% & $+4.3$ \\
R4 & 60--82   & 178 & 2{,}490 & 48.2\,\% & 52.5\,\% & $+4.3$ \\
R5 & 83--121  & 164 & 3{,}296 & 51.1\,\% & 55.4\,\% & $+4.3$ \\
R6 & 122--200 & 170 & 5{,}368 & 59.8\,\% & 64.1\,\% & $+4.3$ \\
\bottomrule
\end{tabular}

\vspace{3pt}
\begin{minipage}{\textwidth}\footnotesize
Baselines, group sizes and prediction counts are measured on the test file. The
model column carries the run's overall $+4.3\pt$ gap across to each group
rather than reporting a per-group measurement, and should be read as such: the
column that is measured here is the baseline, and what it shows is that the
difficulty of a user is set before any model is involved.
\end{minipage}
\end{table}
